%% ============================================================
%%  Chapter 6 — Exception Handling
%%  06_exception_handling.tex
%% ============================================================
\chapter{Exception Handling}
\label{chap:exceptions}

\chapterintro{%
  This chapter documents the custom exception architecture of the Clinic
  Management API. It covers the design philosophy behind using typed
  exceptions as the primary error-signalling mechanism, provides the
  full PHP source for every custom exception class, explains how
  exceptions are registered and rendered in the Laravel 12 bootstrap
  pipeline, and concludes with a complete mapping of every domain error
  condition to its canonical HTTP status code and JSON response envelope.
}

%% ─────────────────────────────────────────────────────────────────────────
\section{Exception Architecture Philosophy}
\label{sec:exception-philosophy}
%% ─────────────────────────────────────────────────────────────────────────

A naive API implementation handles errors by scattering conditional
\texttt{if} statements and \texttt{return response()->json([...], 4xx)}
calls throughout controllers and service classes. This approach has
several compounding problems: the same error shape must be reproduced at
every call site, the HTTP status code is chosen ad-hoc rather than
systematically, and there is no single location where a developer can
audit the full set of error conditions the system can produce.

The Clinic Management API takes a different approach. Every domain error
condition is represented as a \textbf{named, typed PHP exception class}.
Service classes \texttt{throw} these exceptions to signal that a business
rule has been violated. Controllers contain no error-handling logic at
all. The Laravel exception handler --- registered once in
\texttt{bootstrap/app.php} --- intercepts every exception, matches it
against the registered type map, and renders the canonical JSON response.

This architecture has three material consequences:

\begin{enumerate}
  \item \textbf{Single source of truth.} The HTTP status code and error
    message for any given domain error are defined in exactly one place:
    the exception class's \texttt{render()} method. They cannot drift
    between call sites.

  \item \textbf{Testability.} PHPUnit feature tests can assert on the
    HTTP status code and JSON body without knowing anything about the
    internal implementation. They simply trigger the business condition
    and observe the response.

  \item \textbf{Service-layer purity.} Service classes express domain
    intent --- \textit{``this appointment slot is already taken''} ---
    rather than HTTP concerns. The translation from domain language to
    HTTP protocol is the exception handler's responsibility, not the
    service's.
\end{enumerate}

\begin{figure}[h!]
  \centering
  \begin{tcolorbox}[
    enhanced, arc=4pt, boxrule=1pt,
    colback=LightGray, colframe=RuleGray,
    left=12pt, right=12pt, top=10pt, bottom=10pt,
    width=\linewidth,
  ]
  \centering
  \textcolor{PrimaryBlue}{\textbf{Service Layer}}
  \quad\textcolor{red!60!black}{$\xrightarrow{\texttt{throw}}$}\quad
  \textcolor{MethodDelete}{\textbf{Typed Exception}}
  \quad\textcolor{AccentTeal}{$\xrightarrow{\texttt{render()}}$}\quad
  \textcolor{MethodPost}{\textbf{JSON Response}}

  \vspace{6pt}
  \small
  Controllers never catch exceptions. The Laravel exception handler
  pipeline calls \texttt{render()} on every unhandled throwable.
  \end{tcolorbox}
  \caption{The exception flow from service layer to JSON HTTP response.}
  \label{fig:exception-flow}
\end{figure}

%% ─────────────────────────────────────────────────────────────────────────
\section{Exception Class Hierarchy}
\label{sec:exception-hierarchy}
%% ─────────────────────────────────────────────────────────────────────────

All custom exception classes reside in the
\texttt{app/Exceptions/} directory. They extend PHP's built-in
\texttt{Exception} class directly. Each class implements a
\texttt{render()} method that returns an
\texttt{Illuminate\textbackslash{}Http\textbackslash{}JsonResponse}
with the correct HTTP status code and the standard API envelope.

\begin{table}[h!]
  \centering
  \caption{Full inventory of custom exception classes, their HTTP status
    codes, and the domain conditions that trigger them.}
  \label{tab:exception-inventory}
  \renewcommand{\arraystretch}{1.5}
  \begin{tabularx}{\linewidth}{@{}L{5.2cm} C{2.0cm} X@{}}
    \toprule
    \rowcolor{PrimaryBlue}
    \tblhdr{Exception Class} & \tblhdr{HTTP Status} & \tblhdr{Trigger Condition} \\
    \midrule
    \phpclass{AccountDeactivatedException}
      & \stforbid{}
      & Authenticated user's \texttt{is\_active} flag is \texttt{false}
        at login time. \\
    \rowalt
    \phpclass{AppointmentConflictException}
      & \stconflict{}
      & A booking or reschedule request targets a time slot that is
        already occupied by a non-cancelled appointment. \\
    \phpclass{ClinicScopeViolationException}
      & \stforbid{}
      & A user attempts to read or write a resource belonging to a
        different clinic, or a doctor attempts to access another
        doctor's prescription. \\
    \rowalt
    \phpclass{InvalidAppointmentStateException}
      & \stunproc{}
      & An operation (reschedule, prescription creation, status
        transition) is attempted on an appointment whose current
        \texttt{status} value does not permit it. \\
    \phpclass{SelfDemotionException}
      & \stforbid{}
      & A Super Administrator attempts to change their own role via
        the \route{PATCH /api/v1/admin/users/\{id\}/role} endpoint. \\
    \rowalt
    \phpclass{RoleUnauthorizedException}
      & \stforbid{}
      & Programmatic role check failure when the
        \phpclass{RoleMiddleware} is bypassed or additional in-service
        role validation is required. \\
    \phpclass{ResourceNotFoundException}
      & \stnotfound{}
      & Manual not-found signal raised by a service or repository when
        route model binding is not in use (e.g.\ nested resource
        lookup). \\
    \rowalt
    \phpclass{DuplicateScheduleException}
      & \stunproc{}
      & A doctor attempts to create a second schedule entry for a day
        of the week that already has one (\texttt{day\_of\_week}
        uniqueness violation caught before DB constraint fires). \\
    \bottomrule
  \end{tabularx}
\end{table}

%% ─────────────────────────────────────────────────────────────────────────
\section{The ApiResponse Trait}
\label{sec:apiresponse-trait}
%% ─────────────────────────────────────────────────────────────────────────

Before examining individual exception classes it is useful to understand
the \phpclass{ApiResponse} trait, which is used by both exception
\texttt{render()} methods and controllers to guarantee a consistent
JSON envelope shape across every response the API produces.

\begin{lstlisting}[style=php, caption={\texttt{app/Traits/ApiResponse.php}
  --- Centralised response factory used by exceptions and controllers.}]
<?php

namespace App\Traits;

use Illuminate\Http\JsonResponse;

trait ApiResponse
{
    protected function successResponse(
        mixed  $data,
        string $message = 'Operation completed successfully.',
        int    $status  = 200,
    ): JsonResponse {
        return response()->json([
            'success' => true,
            'message' => $message,
            'data'    => $data,
        ], $status);
    }

    protected function errorResponse(
        string $message,
        int    $status,
        mixed  $errors = null,
    ): JsonResponse {
        $body = [
            'success' => false,
            'message' => $message,
        ];

        if ($errors !== null) {
            $body['errors'] = $errors;
        }

        return response()->json($body, $status);
    }
}
\end{lstlisting}

Exception classes call the standalone \texttt{response()->json()} helper
directly (since they do not have access to the \phpclass{ApiResponse}
trait) but produce structurally identical envelopes.

%% ─────────────────────────────────────────────────────────────────────────
\section{Custom Exception Implementations}
\label{sec:exception-implementations}
%% ─────────────────────────────────────────────────────────────────────────

%% ── §4.1 ─────────────────────────────────────────────────────────────────
\subsection{AccountDeactivatedException}
\label{sec:ex-account-deactivated}

Thrown by \phpclass{AuthService} when a user who supplies valid
credentials has \texttt{is\_active = false}. The service deliberately
verifies credentials first and checks the active flag second to prevent
user-enumeration attacks.

\begin{lstlisting}[style=php, caption={\texttt{app/Exceptions/AccountDeactivatedException.php}.}]
<?php

namespace App\Exceptions;

use Exception;
use Illuminate\Http\JsonResponse;

class AccountDeactivatedException extends Exception
{
    public function __construct()
    {
        parent::__construct(
            'Your account has been deactivated. Please contact your administrator.'
        );
    }

    public function render(): JsonResponse
    {
        return response()->json([
            'success' => false,
            'message' => $this->getMessage(),
        ], 403);
    }
}
\end{lstlisting}

%% ── §4.2 ─────────────────────────────────────────────────────────────────
\subsection{AppointmentConflictException}
\label{sec:ex-appointment-conflict}

Thrown by \phpclass{AppointmentService} when the requested time slot
is already occupied by a non-cancelled appointment. It is raised in
two situations: when the distributed cache lock cannot be acquired
(indicating a concurrent request is mid-write) and when the slot
availability check determines the slot is booked.

\begin{lstlisting}[style=php, caption={\texttt{app/Exceptions/AppointmentConflictException.php}.}]
<?php

namespace App\Exceptions;

use Exception;
use Illuminate\Http\JsonResponse;

class AppointmentConflictException extends Exception
{
    public function __construct(
        string $message = 'The requested appointment slot is not available.'
    ) {
        parent::__construct($message);
    }

    public function render(): JsonResponse
    {
        return response()->json([
            'success' => false,
            'message' => $this->getMessage(),
        ], 409);
    }
}
\end{lstlisting}

\begin{notebox}[Why 409 Conflict?]
  RFC~7231 defines HTTP~409 Conflict as appropriate when ``the request
  could not be completed due to a conflict with the current state of
  the target resource.'' A double-booking is precisely such a conflict:
  the slot resource is in a state (occupied) that is incompatible with
  the operation (booking). HTTP~422 would be semantically incorrect here
  because the request payload is valid --- the slot simply is not free.
\end{notebox}

%% ── §4.3 ─────────────────────────────────────────────────────────────────
\subsection{ClinicScopeViolationException}
\label{sec:ex-clinic-scope}

Thrown by any controller or service that detects a cross-clinic resource
access attempt. This exception serves two distinct security purposes:

\begin{enumerate}
  \item It prevents a user from reading or writing resources belonging to
    a different clinic (horizontal privilege escalation).
  \item It prevents a Doctor from accessing another Doctor's prescription
    within the same clinic (intra-role horizontal escalation).
\end{enumerate}

The exception always renders as \stforbid{} with a deliberately generic
message that does not confirm whether the target resource exists.

\begin{lstlisting}[style=php, caption={\texttt{app/Exceptions/ClinicScopeViolationException.php}.}]
<?php

namespace App\Exceptions;

use Exception;
use Illuminate\Http\JsonResponse;

class ClinicScopeViolationException extends Exception
{
    public function __construct(
        string $message = 'You do not have permission to access this resource.'
    ) {
        parent::__construct($message);
    }

    public function render(): JsonResponse
    {
        return response()->json([
            'success' => false,
            'message' => $this->getMessage(),
        ], 403);
    }
}
\end{lstlisting}

\begin{warningbox}[Indistinguishable from Role Mismatch]
  The HTTP~403 response produced by \phpclass{ClinicScopeViolationException}
  is intentionally identical in structure to the response produced by
  \phpclass{RoleMiddleware} for a role mismatch. An attacker who
  correctly guesses a resource UUID from a different clinic cannot
  distinguish ``you lack the role'' from ``the resource exists but is
  not yours.'' This eliminates oracle-style UUID enumeration via
  differential error analysis.
\end{warningbox}

%% ── §4.4 ─────────────────────────────────────────────────────────────────
\subsection{InvalidAppointmentStateException}
\label{sec:ex-invalid-state}

Thrown whenever an operation is attempted on an appointment whose
current \texttt{status} value renders the operation impossible. The
four concrete scenarios are:

\begin{itemize}
  \item Rescheduling a \texttt{completed} appointment.
  \item Rescheduling a \texttt{cancelled} appointment.
  \item Creating a prescription for a \texttt{cancelled} appointment.
  \item Attempting an invalid status transition (e.g.\
    \texttt{completed} $\to$ \texttt{confirmed}).
\end{itemize}

The exception accepts a contextual message at the call site so that the
error message describes the specific violation.

\begin{lstlisting}[style=php, caption={\texttt{app/Exceptions/InvalidAppointmentStateException.php}.}]
<?php

namespace App\Exceptions;

use Exception;
use Illuminate\Http\JsonResponse;

class InvalidAppointmentStateException extends Exception
{
    public function __construct(
        string $message = 'This operation cannot be performed on the appointment in its current state.'
    ) {
        parent::__construct($message);
    }

    public function render(): JsonResponse
    {
        return response()->json([
            'success' => false,
            'message' => $this->getMessage(),
        ], 422);
    }
}
\end{lstlisting}

Example usage in \phpclass{AppointmentService}:

\begin{lstlisting}[style=php, caption={Contextual \phpclass{InvalidAppointmentStateException}
  throws within \texttt{AppointmentService}.}]
public function reschedule(Appointment $appointment, array $data): Appointment
{
    if ($appointment->status === AppointmentStatus::COMPLETED) {
        throw new InvalidAppointmentStateException(
            'A completed appointment cannot be rescheduled.'
        );
    }

    if ($appointment->status === AppointmentStatus::CANCELLED) {
        throw new InvalidAppointmentStateException(
            'A cancelled appointment cannot be rescheduled.'
        );
    }

    // Slot availability check and update ...
}

public function updateStatus(Appointment $appointment, AppointmentStatus $newStatus): Appointment
{
    $allowed = match ($appointment->status) {
        AppointmentStatus::PENDING    => [AppointmentStatus::CONFIRMED],
        AppointmentStatus::CONFIRMED  => [AppointmentStatus::COMPLETED],
        default                       => [],
    };

    if (! in_array($newStatus, $allowed, strict: true)) {
        throw new InvalidAppointmentStateException(
            "Cannot transition appointment from {$appointment->status->value} to {$newStatus->value}."
        );
    }

    return $this->appointmentRepository->updateStatus($appointment, $newStatus);
}
\end{lstlisting}

%% ── §4.5 ─────────────────────────────────────────────────────────────────
\subsection{SelfDemotionException}
\label{sec:ex-self-demotion}

Thrown exclusively by \phpclass{UserController::updateRole()} when the
authenticated Super Administrator attempts to change their own role. This
is a hard-coded guard that cannot be bypassed by any configuration or
environment variable.

\begin{lstlisting}[style=php, caption={\texttt{app/Exceptions/SelfDemotionException.php}.}]
<?php

namespace App\Exceptions;

use Exception;
use Illuminate\Http\JsonResponse;

class SelfDemotionException extends Exception
{
    public function __construct()
    {
        parent::__construct('You cannot change your own role.');
    }

    public function render(): JsonResponse
    {
        return response()->json([
            'success' => false,
            'message' => $this->getMessage(),
        ], 403);
    }
}
\end{lstlisting}

%% ── §4.6 ─────────────────────────────────────────────────────────────────
\subsection{RoleUnauthorizedException}
\label{sec:ex-role-unauthorized}

While role enforcement is primarily handled by \phpclass{RoleMiddleware}
at the route group level, there are scenarios in which a service class
must perform a secondary role check that cannot be expressed as a route
group constraint --- for example, verifying that the doctor initiating a
prescription is the same doctor assigned to the appointment.
\phpclass{RoleUnauthorizedException} provides a typed exception for
these in-service role assertions.

\begin{lstlisting}[style=php, caption={\texttt{app/Exceptions/RoleUnauthorizedException.php}.}]
<?php

namespace App\Exceptions;

use Exception;
use Illuminate\Http\JsonResponse;

class RoleUnauthorizedException extends Exception
{
    public function __construct(
        string $message = 'You do not have permission to perform this action.'
    ) {
        parent::__construct($message);
    }

    public function render(): JsonResponse
    {
        return response()->json([
            'success' => false,
            'message' => $this->getMessage(),
        ], 403);
    }
}
\end{lstlisting}

%% ── §4.7 ─────────────────────────────────────────────────────────────────
\subsection{ResourceNotFoundException}
\label{sec:ex-not-found}

Laravel's route model binding automatically throws a
\texttt{ModelNotFoundException} when a UUID path parameter cannot be
resolved to an Eloquent record, and this is mapped to \stnotfound{}
in the bootstrap handler. However, service classes that perform manual
lookups (e.g.\ resolving a nested resource by composite key rather than
a single UUID) use \phpclass{ResourceNotFoundException} to produce the
same standardised response.

\begin{lstlisting}[style=php, caption={\texttt{app/Exceptions/ResourceNotFoundException.php}.}]
<?php

namespace App\Exceptions;

use Exception;
use Illuminate\Http\JsonResponse;

class ResourceNotFoundException extends Exception
{
    public function __construct(
        string $resource = 'Resource',
        string $id       = '',
    ) {
        $suffix  = $id !== '' ? " with identifier [{$id}]" : '';
        parent::__construct("{$resource}{$suffix} was not found.");
    }

    public function render(): JsonResponse
    {
        return response()->json([
            'success' => false,
            'message' => $this->getMessage(),
        ], 404);
    }
}
\end{lstlisting}

Example usage:

\begin{lstlisting}[style=php, caption={Manual not-found signal in a repository.}]
public function findByCompositeKey(string $doctorId, int $dayOfWeek): DoctorSchedule
{
    $schedule = DoctorSchedule::where('doctor_id',    $doctorId)
                               ->where('day_of_week', $dayOfWeek)
                               ->first();

    if ($schedule === null) {
        throw new ResourceNotFoundException('DoctorSchedule');
    }

    return $schedule;
}
\end{lstlisting}

%% ── §4.8 ─────────────────────────────────────────────────────────────────
\subsection{DuplicateScheduleException}
\label{sec:ex-duplicate-schedule}

Thrown by \phpclass{DoctorScheduleService} when a doctor attempts to
create a schedule entry for a day of the week that already has one.
The service catches this condition before the database constraint fires,
enabling a user-friendly message rather than a raw PDO exception.

\begin{lstlisting}[style=php, caption={\texttt{app/Exceptions/DuplicateScheduleException.php}.}]
<?php

namespace App\Exceptions;

use Exception;
use Illuminate\Http\JsonResponse;

class DuplicateScheduleException extends Exception
{
    public function __construct(int $dayOfWeek)
    {
        $dayName = match ($dayOfWeek) {
            0 => 'Sunday',
            1 => 'Monday',
            2 => 'Tuesday',
            3 => 'Wednesday',
            4 => 'Thursday',
            5 => 'Friday',
            6 => 'Saturday',
            default => "day {$dayOfWeek}",
        };

        parent::__construct(
            "A schedule entry for {$dayName} already exists for this doctor."
        );
    }

    public function render(): JsonResponse
    {
        return response()->json([
            'success' => false,
            'message' => $this->getMessage(),
        ], 422);
    }
}
\end{lstlisting}

%% ─────────────────────────────────────────────────────────────────────────
\section{Exception Registration in the Bootstrap Pipeline}
\label{sec:exception-bootstrap}
%% ─────────────────────────────────────────────────────────────────────────

In Laravel 12, the application is bootstrapped entirely through
\texttt{bootstrap/app.php}. There is no \texttt{App\textbackslash{}Exceptions\textbackslash{}Handler}
class; instead, exception rendering and reporting behaviour is
configured through the \texttt{withExceptions()} callback passed to
\texttt{Application::configure()}.

\subsection{The withExceptions Callback}

The \texttt{withExceptions()} callback receives an
\texttt{Illuminate\textbackslash{}Foundation\textbackslash{}Configuration\textbackslash{}Exceptions}
instance. Two methods are available:

\begin{itemize}
  \item \texttt{render(callable \$callback)} --- registers a rendering
    closure that intercepts a specific exception type and returns a
    \texttt{Response}. If the closure returns \texttt{null}, Laravel
    falls through to the next renderer.
  \item \texttt{report(callable \$callback)} --- registers a reporting
    closure that logs or ships the exception to an error-tracking
    service (e.g.\ Sentry, Bugsnag) without affecting the HTTP response.
\end{itemize}

\subsection{Full Bootstrap Configuration}

\begin{lstlisting}[style=php, caption={\texttt{bootstrap/app.php} ---
  Complete exception handler registration for all custom exception
  types and Laravel framework exceptions.}]
<?php

use App\Exceptions\AccountDeactivatedException;
use App\Exceptions\AppointmentConflictException;
use App\Exceptions\ClinicScopeViolationException;
use App\Exceptions\DuplicateScheduleException;
use App\Exceptions\InvalidAppointmentStateException;
use App\Exceptions\ResourceNotFoundException;
use App\Exceptions\RoleUnauthorizedException;
use App\Exceptions\SelfDemotionException;
use App\Http\Middleware\ClinicScopeMiddleware;
use App\Http\Middleware\RoleMiddleware;
use Illuminate\Auth\AuthenticationException;
use Illuminate\Database\Eloquent\ModelNotFoundException;
use Illuminate\Foundation\Application;
use Illuminate\Foundation\Configuration\Exceptions;
use Illuminate\Foundation\Configuration\Middleware;
use Illuminate\Validation\ValidationException;
use Symfony\Component\HttpKernel\Exception\MethodNotAllowedHttpException;
use Symfony\Component\HttpKernel\Exception\NotFoundHttpException;

return Application::configure(basePath: dirname(__DIR__))
    ->withRouting(
        api:      __DIR__.'/../routes/api.php',
        apiPrefix: 'api',
        commands: __DIR__.'/../routes/console.php',
        health:   '/up',
    )
    ->withMiddleware(function (Middleware $middleware) {
        $middleware->alias([
            'role'         => RoleMiddleware::class,
            'clinic.scope' => ClinicScopeMiddleware::class,
        ]);
    })
    ->withExceptions(function (Exceptions $exceptions) {

        /*
        |------------------------------------------------------------------
        | Framework Exceptions
        |------------------------------------------------------------------
        | Laravel and Symfony throw several built-in exceptions that must
        | be intercepted and shaped into the standard API envelope.
        |------------------------------------------------------------------
        */

        // 401 — bearer token missing or revoked
        $exceptions->render(function (AuthenticationException $e) {
            return response()->json([
                'success' => false,
                'message' => 'Unauthenticated.',
            ], 401);
        });

        // 404 — route model binding failure
        $exceptions->render(function (ModelNotFoundException $e) {
            return response()->json([
                'success' => false,
                'message' => 'Resource not found.',
            ], 404);
        });

        // 404 — unregistered route
        $exceptions->render(function (NotFoundHttpException $e) {
            return response()->json([
                'success' => false,
                'message' => 'The requested endpoint does not exist.',
            ], 404);
        });

        // 405 — wrong HTTP method for a registered route
        $exceptions->render(function (MethodNotAllowedHttpException $e) {
            return response()->json([
                'success' => false,
                'message' => 'HTTP method not allowed for this endpoint.',
            ], 405);
        });

        // 422 — form request validation failure
        $exceptions->render(function (ValidationException $e) {
            return response()->json([
                'success' => false,
                'errors'  => $e->errors(),
            ], 422);
        });

        /*
        |------------------------------------------------------------------
        | Domain Exceptions
        |------------------------------------------------------------------
        | All custom exceptions implement render() directly, so Laravel
        | will call render() on them automatically. The registrations
        | below are provided for explicit documentation and to allow
        | report() hooks to be attached to specific types.
        |------------------------------------------------------------------
        */

        // 403 — account deactivated after valid credentials
        $exceptions->render(function (AccountDeactivatedException $e) {
            return $e->render();
        });

        // 409 — appointment slot conflict
        $exceptions->render(function (AppointmentConflictException $e) {
            return $e->render();
        });

        // 403 — cross-clinic resource access
        $exceptions->render(function (ClinicScopeViolationException $e) {
            return $e->render();
        });

        // 422 — duplicate doctor schedule day
        $exceptions->render(function (DuplicateScheduleException $e) {
            return $e->render();
        });

        // 422 — invalid appointment state transition
        $exceptions->render(function (InvalidAppointmentStateException $e) {
            return $e->render();
        });

        // 404 — manual resource not found from service layer
        $exceptions->render(function (ResourceNotFoundException $e) {
            return $e->render();
        });

        // 403 — in-service role assertion failure
        $exceptions->render(function (RoleUnauthorizedException $e) {
            return $e->render();
        });

        // 403 — super admin self-demotion attempt
        $exceptions->render(function (SelfDemotionException $e) {
            return $e->render();
        });

    })->create();
\end{lstlisting}

\begin{notebox}[render() Delegation Pattern]
  Because each custom exception class already implements a \texttt{render()}
  method that returns the correct \texttt{JsonResponse}, the registrations
  in the bootstrap file simply delegate back to \texttt{\$e->render()}.
  This is by design: the HTTP status code and message are owned by the
  exception class, not by the bootstrap file. The bootstrap registrations
  exist to make the exception inventory explicit and to provide
  attachment points for \texttt{report()} hooks without changing the
  render logic.
\end{notebox}

\subsection{Automatic render() Detection}

Laravel's exception handler checks whether an unhandled throwable
implements a \texttt{render()} method before consulting the registered
closures. This means that even without the explicit \texttt{render}
registrations in the bootstrap file, custom exception classes with a
\texttt{render()} method would produce the correct JSON responses. The
explicit registrations are included for documentation and
\texttt{report()} hook purposes, not for functional necessity.

%% ─────────────────────────────────────────────────────────────────────────
\section{The AppointmentStatus Enum and State Machine}
\label{sec:appointment-state-machine}
%% ─────────────────────────────────────────────────────────────────────────

The \phpclass{InvalidAppointmentStateException} is the enforcement
mechanism for the appointment state machine. Understanding the enum and
the transition guard is essential to understanding when and why this
exception fires.

\begin{lstlisting}[style=php, caption={\texttt{app/Enums/AppointmentStatus.php}
  --- Backed string enum encoding all valid appointment states.}]
<?php

namespace App\Enums;

enum AppointmentStatus: string
{
    case PENDING   = 'pending';
    case CONFIRMED = 'confirmed';
    case COMPLETED = 'completed';
    case CANCELLED = 'cancelled';

    public function isTerminal(): bool
    {
        return match ($this) {
            self::COMPLETED,
            self::CANCELLED => true,
            default         => false,
        };
    }

    public function allowedTransitions(): array
    {
        return match ($this) {
            self::PENDING   => [self::CONFIRMED, self::CANCELLED],
            self::CONFIRMED => [self::COMPLETED, self::CANCELLED],
            self::COMPLETED,
            self::CANCELLED => [],
        };
    }

    public function canTransitionTo(self $target): bool
    {
        return in_array($target, $this->allowedTransitions(), strict: true);
    }
}
\end{lstlisting}

The \texttt{canTransitionTo()} method encapsulates the entire transition
matrix. The service layer calls it and throws
\phpclass{InvalidAppointmentStateException} if it returns \texttt{false}:

\begin{lstlisting}[style=php, caption={State transition guard in
  \texttt{AppointmentService::updateStatus()}.}]
public function updateStatus(
    Appointment       $appointment,
    AppointmentStatus $newStatus,
): Appointment {
    if (! $appointment->status->canTransitionTo($newStatus)) {
        throw new InvalidAppointmentStateException(
            "Cannot transition from [{$appointment->status->value}]"
            . " to [{$newStatus->value}]."
        );
    }

    return $this->appointmentRepository->updateStatus($appointment, $newStatus);
}
\end{lstlisting}

%% ─────────────────────────────────────────────────────────────────────────
\section{Handling Laravel Framework Exceptions}
\label{sec:framework-exceptions}
%% ─────────────────────────────────────────────────────────────────────────

Several exceptions thrown by the Laravel framework itself must also be
intercepted and shaped into the standard JSON envelope. Without
interception these exceptions would produce HTML error pages (in
non-production environments) or empty responses (in production), both
of which are incompatible with a JSON API contract.

\subsection{AuthenticationException}

Thrown by the \texttt{auth:api} Passport middleware when no valid
bearer token is present. Without interception Laravel returns a redirect
response (appropriate for web sessions but not for an API).

\begin{lstlisting}[style=json, caption={Rendered HTTP 401 response.}]
{
  "success": false,
  "message": "Unauthenticated."
}
\end{lstlisting}

\subsection{ModelNotFoundException}

Thrown by Laravel's route model binding when a UUID path parameter
resolves to no Eloquent record. The binding calls
\texttt{findOrFail()}, which throws \texttt{ModelNotFoundException} if
the record does not exist.

\begin{lstlisting}[style=json, caption={Rendered HTTP 404 response.}]
{
  "success": false,
  "message": "Resource not found."
}
\end{lstlisting}

\subsection{ValidationException}

Thrown by \phpclass{FormRequest} subclasses when their
\texttt{rules()} constraints are not satisfied. Laravel normally
redirects back with errors for web routes; for API routes it produces
a raw JSON response that lacks the \texttt{success} key. The
registered renderer reshapes it into the standard envelope.

\begin{lstlisting}[style=json, caption={Rendered HTTP 422 response for a
  failed login validation.}]
{
  "success": false,
  "errors": {
    "email":    ["The email field is required."],
    "password": ["The password field is required."]
  }
}
\end{lstlisting}

\subsection{NotFoundHttpException}

Thrown by the router when a request arrives for a URL that does not
match any registered route. Distinct from \texttt{ModelNotFoundException},
which concerns resource lookup within a matched route.

\begin{lstlisting}[style=json, caption={Rendered HTTP 404 response for
  an unregistered route.}]
{
  "success": false,
  "message": "The requested endpoint does not exist."
}
\end{lstlisting}

\subsection{MethodNotAllowedHttpException}

Thrown by the router when the HTTP method used does not match the
registered methods for a given path (e.g.\ a \texttt{GET} request to
a \texttt{POST}-only endpoint).

\begin{lstlisting}[style=json, caption={Rendered HTTP 405 response.}]
{
  "success": false,
  "message": "HTTP method not allowed for this endpoint."
}
\end{lstlisting}

%% ─────────────────────────────────────────────────────────────────────────
\section{Domain Error to HTTP Status Code Mapping}
\label{sec:error-status-map}
%% ─────────────────────────────────────────────────────────────────────────

The following table is the definitive reference for every error condition
the API can produce, the exception or mechanism that generates it, its
HTTP status code, and the JSON envelope shape.

\renewcommand{\arraystretch}{1.45}

\begin{longtable}{@{}C{2.0cm} L{4.8cm} X@{}}
  \toprule
  \rowcolor{PrimaryBlue}
  \tblhdr{HTTP Status} &
  \tblhdr{Source / Exception} &
  \tblhdr{Trigger Condition and Response Shape} \\
  \midrule
  \endfirsthead

  \toprule
  \rowcolor{PrimaryBlue}
  \tblhdr{HTTP Status} &
  \tblhdr{Source / Exception} &
  \tblhdr{Trigger Condition and Response Shape (cont.)} \\
  \midrule
  \endhead

  \bottomrule
  \endfoot

  \multicolumn{3}{@{}l}{\small\textit{--- 400-range client errors ---}} \\[4pt]

  \stunauth{}
    & \texttt{AuthenticationException}
      \newline (Laravel framework)
    & Bearer token absent, malformed, or revoked. The \texttt{auth:api}
      Passport middleware rejects the request before it reaches the
      application.
      \newline\texttt{\{"success":false,"message":"Unauthenticated."\}} \\
  \rowalt

  \stforbid{}
    & \phpclass{RoleMiddleware}
    & Authenticated user's role is not in the list of roles allowed for
      this route group.
      \newline\texttt{\{"success":false,"message":"You do not have permission to perform this action."\}} \\

  \stforbid{}
    & \phpclass{AccountDeactivatedException}
    & Valid credentials supplied at login, but \texttt{is\_active = false}.
      \newline\texttt{\{"success":false,"message":"Your account has been deactivated..."\}} \\
  \rowalt

  \stforbid{}
    & \phpclass{ClinicScopeViolationException}
    & UUID resource belongs to a different clinic, or prescription belongs
      to a different doctor.
      \newline\texttt{\{"success":false,"message":"You do not have permission to access this resource."\}} \\

  \stforbid{}
    & \phpclass{SelfDemotionException}
    & Super Administrator calls \texttt{PATCH .../users/\{id\}/role}
      where \texttt{\{id\}} is their own UUID.
      \newline\texttt{\{"success":false,"message":"You cannot change your own role."\}} \\
  \rowalt

  \stforbid{}
    & \phpclass{RoleUnauthorizedException}
    & Programmatic in-service role assertion fails (secondary ownership
      check beyond route-group middleware).
      \newline\texttt{\{"success":false,"message":"You do not have permission to perform this action."\}} \\

  \stnotfound{}
    & \texttt{ModelNotFoundException}
      \newline (Laravel framework)
    & Route model binding cannot resolve a UUID path parameter to an
      existing Eloquent record.
      \newline\texttt{\{"success":false,"message":"Resource not found."\}} \\
  \rowalt

  \stnotfound{}
    & \texttt{NotFoundHttpException}
      \newline (Symfony/Laravel)
    & Request URL does not match any registered route.
      \newline\texttt{\{"success":false,"message":"The requested endpoint does not exist."\}} \\

  \stnotfound{}
    & \phpclass{ResourceNotFoundException}
    & Manual not-found signal from a service or repository performing a
      composite-key lookup outside of route model binding.
      \newline\texttt{\{"success":false,"message":"<Resource> was not found."\}} \\
  \rowalt

  405
    & \texttt{MethodNotAllowedHttpException}
      \newline (Symfony/Laravel)
    & Registered route exists for the path but not for the supplied
      HTTP method.
      \newline\texttt{\{"success":false,"message":"HTTP method not allowed for this endpoint."\}} \\

  \stconflict{}
    & \phpclass{AppointmentConflictException}
    & Booking or reschedule request targets a slot already occupied by a
      non-cancelled appointment, or the distributed cache lock cannot be
      acquired due to a concurrent booking write.
      \newline\texttt{\{"success":false,"message":"The requested appointment slot is not available."\}} \\
  \rowalt

  \stunproc{}
    & \texttt{ValidationException}
      \newline (Laravel framework)
    & \phpclass{FormRequest} validation rules are not satisfied.
      Returns structured field-level errors under the \texttt{errors} key.
      \newline\texttt{\{"success":false,"errors":\{"field":["msg"]\}\}} \\

  \stunproc{}
    & \phpclass{InvalidAppointmentStateException}
    & An operation (reschedule, prescription authoring, status
      transition) is attempted on an appointment in an incompatible
      state.
      \newline\texttt{\{"success":false,"message":"<context-specific message>"\}} \\
  \rowalt

  \stunproc{}
    & \phpclass{DuplicateScheduleException}
    & Doctor attempts to create a second schedule entry for a
      \texttt{day\_of\_week} that already has one.
      \newline\texttt{\{"success":false,"message":"A schedule entry for <Day> already exists for this doctor."\}} \\

\end{longtable}

%% ─────────────────────────────────────────────────────────────────────────
\section{Exception Rendering in the Testing Environment}
\label{sec:exception-testing}
%% ─────────────────────────────────────────────────────────────────────────

PHPUnit feature tests exercise the full HTTP stack including the exception
handler. Laravel's \texttt{withoutExceptionHandling()} helper can be
called in a test to suppress the handler and allow the raw exception to
propagate, which is useful for debugging. The production test suite,
however, leaves the handler active and asserts on the HTTP response
directly:

\begin{lstlisting}[style=php, caption={Feature test asserting the correct
  JSON envelope for a clinic scope violation.}]
public function test_doctor_cannot_access_patient_from_another_clinic(): void
{
    $clinicA = Clinic::factory()->create();
    $clinicB = Clinic::factory()->create();

    $doctor  = User::factory()->doctor()->forClinic($clinicA)->create();
    $patient = Patient::factory()->forClinic($clinicB)->create();

    $response = $this->actingAs($doctor, 'api')
                     ->getJson("/api/v1/doctor/patients/{$patient->id}");

    $response->assertStatus(403)
             ->assertJson([
                 'success' => false,
                 'message' => 'You do not have permission to access this resource.',
             ]);
}
\end{lstlisting}

\begin{lstlisting}[style=php, caption={Feature test asserting the conflict
  response for a double-booking attempt.}]
public function test_booking_an_already_taken_slot_returns_conflict(): void
{
    $clinic    = Clinic::factory()->create();
    $doctor    = User::factory()->doctor()->forClinic($clinic)->create();
    $assistant = User::factory()->assistant()->forClinic($clinic)->create();
    $patientA  = Patient::factory()->forClinic($clinic)->create();
    $patientB  = Patient::factory()->forClinic($clinic)->create();

    DoctorSchedule::factory()->for($doctor)->monday()->create([
        'start_time'    => '09:00',
        'end_time'      => '17:00',
        'slot_duration' => 30,
    ]);

    // First booking succeeds
    Appointment::factory()->create([
        'doctor_id'        => $doctor->id,
        'clinic_id'        => $clinic->id,
        'appointment_date' => '2026-07-07',
        'start_time'       => '09:00',
        'status'           => 'confirmed',
    ]);

    // Second booking for the same slot must be rejected
    $response = $this->actingAs($assistant, 'api')
                     ->postJson('/api/v1/assistant/appointments', [
                         'doctor_id'        => $doctor->id,
                         'patient_id'       => $patientB->id,
                         'appointment_date' => '2026-07-07',
                         'start_time'       => '09:00',
                     ]);

    $response->assertStatus(409)
             ->assertJson([
                 'success' => false,
                 'message' => 'The requested appointment slot is not available.',
             ]);
}
\end{lstlisting}

%% ─────────────────────────────────────────────────────────────────────────
\section{Exception Design Checklist}
\label{sec:exception-checklist}
%% ─────────────────────────────────────────────────────────────────────────

When adding a new domain operation to the API, the following checklist
ensures that error handling is handled consistently.

\begin{table}[h!]
  \centering
  \caption{Checklist for introducing a new exception to the system.}
  \label{tab:exception-checklist}
  \renewcommand{\arraystretch}{1.5}
  \begin{tabularx}{\linewidth}{@{}C{0.6cm} X@{}}
    \toprule
    \rowcolor{PrimaryBlue}
    \tblhdr{\#} & \tblhdr{Step} \\
    \midrule
    1 & Create a new class in \texttt{app/Exceptions/} that extends
        \texttt{Exception}. \\
    \rowalt
    2 & Add a typed constructor that accepts a contextual message
        string with a sensible default. \\
    3 & Implement \texttt{render(): JsonResponse} returning the correct
        HTTP status code and the \texttt{success/message} envelope. \\
    \rowalt
    4 & Register an explicit \texttt{render()} closure in
        \texttt{bootstrap/app.php} for documentation and
        \texttt{report()} hook attachment. \\
    5 & Add the new exception to Table~\ref{tab:exception-inventory}
        in this chapter and to the domain-error map in
        Table~\ref{tab:error-status-map}. \\
    \rowalt
    6 & Write at least one PHPUnit feature test that triggers the
        exception and asserts the expected HTTP status code and JSON
        body. \\
    7 & Ensure the \texttt{throw} site in the service or controller is
        not wrapped in a \texttt{try/catch} that swallows the exception
        before it reaches the handler. \\
    \bottomrule
  \end{tabularx}
\end{table}

\begin{importantbox}[Never Catch in Controllers]
  Controllers must not contain \texttt{try/catch} blocks that catch
  domain exceptions and return manual error responses. Doing so bypasses
  the central exception handler, breaks the single-source-of-truth
  guarantee for error shapes, and makes the exception invisible to any
  registered \texttt{report()} hooks. If a controller needs to respond
  differently to an exception, the correct approach is to register a
  more specific renderer in \texttt{bootstrap/app.php} --- not to
  handle it inline.
\end{importantbox}
